% ============================================================
% CHAPITRE 6 — Introduction à Adobe Photoshop
% ============================================================

\chapter{Introduction à Adobe Photoshop}

\noindent
\textit{Nous voici au début de la deuxième partie de ce cours, consacrée à la maîtrise d'Adobe Photoshop. Dans les chapitres précédents, nous avons posé les fondations théoriques du design graphique : définition, histoire, composition, couleur, typographie, culture visuelle. Ces connaissances sont indispensables, mais elles ne suffisent pas. Il faut maintenant apprendre à les mettre en œuvre avec des outils professionnels.}

\vspace{1em}

\noindent
Adobe Photoshop est l'un des logiciels les plus utilisés dans le monde du design graphique, de la photographie et de la communication visuelle. Il est devenu un standard incontournable, au point que son nom est souvent employé comme un verbe (« photoshoper » une image). Mais Photoshop n'est pas un logiciel magique : c'est un outil puissant qui demande à être apprivoisé. Il ne remplace ni la créativité, ni la culture visuelle, ni la maîtrise des principes du design. Il les prolonge, il les amplifie.

\vspace{1em}

\noindent
Ce chapitre a pour objectif de te faire découvrir Photoshop de manière progressive et structurée. Nous commencerons par une présentation générale du logiciel et de ses domaines d'utilisation. Nous étudierons ensuite les notions fondamentales d'image matricielle, de pixel, de dimension et de résolution. Nous verrons les principaux formats d'image. Puis nous aborderons l'installation et la configuration du logiciel, la découverte de l'interface, la barre de menus et les panneaux, les outils essentiels, la création et l'enregistrement de documents, les raccourcis clavier, et enfin l'organisation de l'espace de travail. À la fin de ce chapitre, tu disposeras des bases nécessaires pour utiliser Photoshop de manière efficace et professionnelle.

% ============================================================
\section{Présentation d'Adobe Photoshop}
% ============================================================

\subsection{Qu'est-ce qu'Adobe Photoshop ?}

\begin{quote}
\textbf{Adobe Photoshop est un logiciel de création et de retouche d'images matricielles, développé par Adobe Inc. Il permet de manipuler des photographies, de créer des compositions graphiques, de réaliser des photomontages, de concevoir des visuels publicitaires, et bien plus encore.}
\end{quote}

Photoshop est un logiciel de \terme{PAO} (Publication Assistée par Ordinateur), plus précisément un logiciel de \terme{traitement d'images matricielles}. Il travaille sur des images composées de pixels, contrairement à Adobe Illustrator qui travaille sur des images vectorielles.

\subsection{Un peu d'histoire}

Photoshop a été créé en \terme{1988} par \terme{Thomas Knoll} et \terme{John Knoll}, deux frères américains. Le logiciel a été racheté par Adobe en 1988 et sa première version commerciale, \terme{Photoshop 1.0}, est sortie en \terme{1990}.

Depuis, Photoshop a connu de nombreuses versions, chacune apportant de nouvelles fonctionnalités :

\begin{itemize}
    \item \terme{1990} : Photoshop 1.0, première version commerciale.
    \item \terme{1994} : Photoshop 3.0, introduction des calques.
    \item \terme{1998} : Photoshop 5.0, introduction de l'historique.
    \item \terme{2003} : Photoshop CS (Creative Suite), intégration à la suite Adobe.
    \item \terme{2013} : Photoshop CC (Creative Cloud), passage à l'abonnement.
    \item \terme{2020} : Photoshop 2020, introduction de l'intelligence artificielle.
    \item \terme{Aujourd'hui} : versions régulières avec nouvelles fonctionnalités (IA générative, etc.).
\end{itemize}

Photoshop est aujourd'hui utilisé par des millions de professionnels et d'amateurs dans le monde entier.

\subsection{Les caractéristiques principales}

\subsubsection{Logiciel matriciel}

Photoshop travaille sur des \terme{images matricielles}, composées de pixels. Cela permet une grande précision dans la retouche, mais limite l'agrandissement sans perte de qualité (voir section 3).

\subsubsection{Logiciel de retouche}

Photoshop excelle dans la \terme{retouche photographique} : correction de l'exposition, des couleurs, des imperfections, retouche beauté, etc.

\subsubsection{Logiciel de composition}

Photoshop permet de \terme{composer} des images : photomontages, collages, compositions publicitaires.

\subsubsection{Logiciel de création}

Photoshop permet de \terme{créer} des visuels : affiches, bannières, mockups, illustrations numériques.

\subsubsection{Logiciel professionnel}

Photoshop est un \terme{standard professionnel}, utilisé dans les agences, les studios, les entreprises. Sa maîtrise est souvent exigée dans les offres d'emploi.

\subsection{Photoshop dans la suite Adobe}

Photoshop fait partie de la \terme{suite Adobe Creative Cloud}, qui comprend également :

\begin{itemize}
    \item \terme{Adobe Illustrator} : création vectorielle.
    \item \terme{Adobe InDesign} : mise en page éditoriale.
    \item \terme{Adobe After Effects} : motion design.
    \item \terme{Adobe Premiere Pro} : montage vidéo.
    \item \terme{Adobe XD / Figma} : design d'interface.
    \item \terme{Adobe Lightroom} : retouche photographique.
\end{itemize}

Ces logiciels sont complémentaires : Photoshop travaille sur les images matricielles, Illustrator sur les images vectorielles, InDesign sur la mise en page. Un designer professionnel maîtrise souvent plusieurs de ces outils.

\subsection{Les versions de Photoshop}

Photoshop existe en plusieurs versions :

\begin{itemize}
    \item \terme{Photoshop CC} : version complète, par abonnement (Creative Cloud).
    \item \terme{Photoshop Elements} : version simplifiée, pour les amateurs.
    \item \terme{Photoshop Express} : version mobile et web, gratuite.
    \item \terme{Photoshop iPad} : version pour tablette.
\end{itemize}

Dans ce cours, nous utiliserons la version \terme{Photoshop CC}, qui est la version professionnelle.

\subsection*{Conclusion de la section}

Adobe Photoshop est un logiciel de création et de retouche d'images matricielles, développé par Adobe. Créé en 1988, il est devenu un standard professionnel. Il permet de retoucher, composer et créer des visuels. Il fait partie de la suite Adobe Creative Cloud, aux côtés d'Illustrator, InDesign, After Effects, etc. La version utilisée dans ce cours est Photoshop CC.

Maintenant que nous avons présenté Photoshop, nous allons, dans la section suivante, étudier ses \terme{domaines d'utilisation}.

% ============================================================
\section{Les domaines d'utilisation de Photoshop}
% ============================================================

\subsection{Pourquoi connaître les domaines d'utilisation ?}

Photoshop est un logiciel très polyvalent, utilisé dans de nombreux domaines. Connaître ces domaines permet de :

\begin{itemize}
    \item \terme{Comprendre} l'étendue des possibilités du logiciel.
    \item \terme{Choisir} les fonctionnalités adaptées à son usage.
    \item \terme{Se positionner} professionnellement.
    \item \terme{Développer} ses compétences dans les domaines pertinents.
\end{itemize}

\subsection{Les principaux domaines d'utilisation}

\subsubsection{La retouche photographique}

C'est le domaine historique de Photoshop. Il s'agit de :

\begin{itemize}
    \item \terme{Corriger} l'exposition, le contraste, la balance des couleurs.
    \item \terme{Supprimer} les imperfections (taches, rides, objets indésirables).
    \item \terme{Améliorer} la netteté, la luminosité, la saturation.
    \item \terme{Retoucher} les portraits (peau, yeux, dents, silhouette).
    \item \terme{Corriger} les dominantes de couleur.
\end{itemize}

Photoshop est utilisé par les photographes professionnels, les studios, les agences de publicité.

\subsubsection{Le photomontage}

Le photomontage consiste à \terme{assembler} plusieurs images pour en créer une nouvelle. Cela implique :

\begin{itemize}
    \item \terme{Détourer} les éléments à assembler.
    \item \terme{Harmoniser} les dimensions, les perspectives, la lumière.
    \item \terme{Créer} des ombres et des reflets réalistes.
    \item \terme{Intégrer} les éléments dans un environnement cohérent.
\end{itemize}

Le photomontage est utilisé dans la publicité, le cinéma, l'illustration.

\subsubsection{La création publicitaire}

Photoshop permet de \terme{créer} des visuels publicitaires :

\begin{itemize}
    \item \terme{Affiches} : annonces, événements, campagnes.
    \item \terme{Bannières} : web, réseaux sociaux.
    \item \terme{Visuels} : produits, packshots.
    \item \terme{Couvertures} : magazines, livres, albums.
\end{itemize}

\subsubsection{Le design web}

Photoshop est utilisé pour créer des \terme{maquettes} de sites web et d'interfaces :

\begin{itemize}
    \item \terme{Maquettes} : pages web, écrans d'application.
    \item \terme{Bannières} : publicités en ligne.
    \item \terme{Visuels} : réseaux sociaux, newsletters.
\end{itemize}

Attention : pour le design web, d'autres outils comme Figma ou Adobe XD sont aujourd'hui plus adaptés. Photoshop reste utile pour créer des visuels statiques.

\subsubsection{Le design d'interface (UI)}

Photoshop a été longtemps utilisé pour concevoir des interfaces d'applications. Aujourd'hui, Figma et Sketch sont plus adaptés, mais Photoshop reste utilisé pour créer des éléments visuels (icônes, illustrations).

\subsubsection{L'illustration numérique}

Photoshop permet de \terme{créer} des illustrations numériques :

\begin{itemize}
    \item \terme{Peinture numérique} : illustration avec pinceaux numériques.
    \item \terme{Concept art} : création de concepts pour le cinéma, le jeu vidéo.
    \item \terme{Illustration éditoriale} : pour la presse, l'édition.
\end{itemize}

\subsubsection{La création de mockups}

Un \terme{mockup} est une maquette réaliste d'un produit (par exemple, un t-shirt, une carte de visite, un packaging). Photoshop est très utilisé pour créer ces mockups, en intégrant des visuels sur des photos de produits.

\subsubsection{Le motion design (limité)}

Photoshop permet de créer des animations simples et des GIF. Pour des animations complexes, on utilise plutôt After Effects. Mais Photoshop reste utile pour préparer des éléments (calques, images) destinés à l'animation.

\subsubsection{La restauration de photos}

Photoshop permet de \terme{restaurer} des photos anciennes ou abîmées :

\begin{itemize}
    \item \terme{Corriger} les rayures, les taches, les déchirures.
    \item \terme{Améliorer} la netteté, le contraste.
    \item \terme{Coloriser} des photos en noir et blanc.
\end{itemize}

\subsubsection{La création de textures et de motifs}

Photoshop permet de créer des \terme{textures} (bois, pierre, tissu, métal) et des \terme{motifs} (répétitions graphiques) utilisés dans d'autres compositions.

\subsection{Tableau récapitulatif des domaines}

\begin{table}[h!]
\centering
\begin{tabular}{@{}p{4cm}p{10cm}@{}}
\toprule
\textbf{Domaine} & \textbf{Utilisation principale} \\
\midrule
Retouche photo & Corriger, améliorer, retoucher des photographies \\
Photomontage & Assembler plusieurs images en une seule \\
Publicité & Créer des affiches, bannières, visuels \\
Design web & Créer des maquettes, bannières, visuels \\
Illustration & Créer des illustrations numériques \\
Mockups & Créer des maquettes réalistes de produits \\
Motion design & Créer des animations simples, des GIF \\
Restauration & Restaurer des photos anciennes \\
Textures & Créer des textures et motifs \\
\bottomrule
\end{tabular}
\caption{Les principaux domaines d'utilisation de Photoshop.}
\end{table}

\subsection{Les limites de Photoshop}

Bien que très polyvalent, Photoshop a des limites :

\begin{itemize}
    \item \terme{Images vectorielles} : Photoshop n'est pas adapté à la création vectorielle (logos, illustrations vectorielles). On utilise Illustrator.
    \item \terme{Mise en page} : Photoshop n'est pas adapté à la mise en page éditoriale (livres, magazines). On utilise InDesign.
    \item \terme{Design d'interface} : Figma ou Sketch sont plus adaptés.
    \item \terme{Motion design} : After Effects est plus adapté.
\end{itemize}

Un designer professionnel utilise le bon outil pour le bon usage.

\subsection*{Conclusion de la section}

Photoshop est utilisé dans de nombreux domaines : retouche photographique, photomontage, création publicitaire, design web, illustration numérique, mockups, motion design, restauration, création de textures. Chaque domaine mobilise des compétences spécifiques. Photoshop a également des limites : il n'est pas adapté aux images vectorielles, à la mise en page éditoriale, au design d'interface, au motion design complexe. Un designer professionnel utilise le bon outil pour le bon usage.

Maintenant que nous connaissons les domaines d'utilisation de Photoshop, nous allons, dans la section suivante, étudier les \terme{images matricielles}.

% ============================================================
\section{Comprendre les images matricielles}
% ============================================================

\subsection{Qu'est-ce qu'une image matricielle ?}

\begin{quote}
\textbf{Une image matricielle (ou bitmap) est une image composée d'une grille de petits carrés appelés pixels. Chaque pixel contient une information de couleur, et l'ensemble des pixels forme l'image.}
\end{quote}

Une image matricielle est donc une \terme{grille de pixels}. Sa qualité dépend de sa \terme{résolution} : plus il y a de pixels, plus l'image est détaillée.

\subsection{Le pixel}

\subsubsection{Définition}

Un \terme{pixel} (\emph{picture element}) est le plus petit élément d'une image matricielle. Il a une position et une couleur.

\subsubsection{Représentation}

Un pixel est souvent représenté comme un petit carré. Mais en réalité, il n'a pas de forme : c'est un point de couleur.

\subsubsection{Composition d'un pixel}

Un pixel contient une information de couleur, codée sur un ou plusieurs canaux :

\begin{itemize}
    \item \terme{Noir et blanc} : 1 canal (niveaux de gris).
    \item \terme{Couleur RVB} : 3 canaux (rouge, vert, bleu).
    \item \terme{Couleur CMJN} : 4 canaux (cyan, magenta, jaune, noir).
    \item \terme{Couleur avec transparence} : 4 canaux (RVB + alpha).
\end{itemize}

\subsection{Avantages et inconvénients des images matricielles}

\subsubsection{Avantages}

\begin{itemize}
    \item \terme{Précision} : permettent une grande précision dans la retouche.
    \item \terme{Richesse} : permettent des dégradés, des nuances, des effets complexes.
    \item \terme{Compatibilité} : formats universels (JPEG, PNG, etc.).
    \item \terme{Photographie} : adaptées à la photographie.
\end{itemize}

\subsubsection{Inconvénients}

\begin{itemize}
    \item \terme{Poids} : fichiers lourds, surtout à haute résolution.
    \item \terme{Perte de qualité à l'agrandissement} : si on agrandit une image matricielle, on perd de la qualité (pixellisation).
    \item \terme{Difficulté à modifier} : on ne peut pas modifier la forme des éléments comme en vectoriel.
\end{itemize}

\subsection{La pixellisation}

\subsubsection{Définition}

La \terme{pixellisation} est l'apparition visible des pixels lorsqu'on agrandit une image matricielle au-delà de sa résolution native.

\subsubsection{Exemple}

Une image de 100 × 100 pixels, agrandie à 1000 × 1000 pixels, montre des pixels visibles : les carrés deviennent apparents, les bords deviennent crénelés.

\subsubsection{Conséquence}

Il est donc essentiel de travailler à la \terme{bonne résolution} dès le départ. On ne peut pas agrandir indéfiniment une image matricielle sans perte de qualité.

\subsection{Comparaison avec les images vectorielles}

\begin{table}[h!]
\centering
\begin{tabular}{@{}p{4cm}p{5cm}p{5cm}@{}}
\toprule
\textbf{Critère} & \textbf{Image matricielle} & \textbf{Image vectorielle} \\
\midrule
Composition & Grille de pixels & Formes géométriques \\
Qualité à l'agrandissement & Perte de qualité & Aucune perte \\
Poids & Souvent lourd & Souvent léger \\
Retouche & Adaptée & Non adaptée \\
Création de logos & Non adaptée & Adaptée \\
Logiciel & Photoshop & Illustrator \\
\bottomrule
\end{tabular}
\caption{Comparaison image matricielle / image vectorielle.}
\end{table}

\subsection{Quand utiliser une image matricielle ?}

\begin{itemize}
    \item \terme{Photographie} : retouche, composition.
    \item \terme{Photomontage} : assemblage d'images.
    \item \terme{Illustration numérique} : peinture numérique.
    \item \terme{Création publicitaire} : affiches, bannières.
    \item \terme{Textures} : création de textures réalistes.
\end{itemize}

\subsection*{Conclusion de la section}

Une \terme{image matricielle} est composée d'une grille de pixels. Chaque pixel contient une information de couleur. Les images matricielles offrent une grande précision mais perdent de la qualité à l'agrandissement (pixellisation). Elles sont adaptées à la photographie, au photomontage, à l'illustration numérique. Elles se distinguent des images vectorielles, qui sont composées de formes géométriques et ne perdent pas de qualité à l'agrandissement.

Maintenant que nous comprenons les images matricielles, nous allons, dans la section suivante, étudier les notions de \terme{pixels, dimensions et résolution}.

% ============================================================
\section{Pixels, dimensions et résolution}
% ============================================================

\subsection{Le pixel}

Nous avons déjà défini le pixel comme le plus petit élément d'une image matricielle. Précisons ici son rôle dans la définition des dimensions et de la résolution.

\subsection{Les dimensions}

\subsubsection{Définition}

Les \terme{dimensions} d'une image sont sa \terme{largeur} et sa \terme{hauteur}, exprimées en pixels, en centimètres, en millimètres, en pouces.

\subsubsection{Modes d'expression}

\begin{itemize}
    \item \terme{En pixels} : 1920 × 1080 px (résolution d'un écran Full HD).
    \item \terme{En centimètres} : 21 × 29,7 cm (format A4).
    \item \terme{En millimètres} : 210 × 297 mm (format A4).
    \item \terme{En pouces} : 8,27 × 11,69 in (format A4).
\end{itemize}

\subsubsection{Choisir les dimensions}

Le choix des dimensions dépend de l'usage :

\begin{itemize}
    \item \terme{Écran} : dimensions en pixels (par exemple, 1920 × 1080 px).
    \item \terme{Impression} : dimensions en centimètres ou millimètres (par exemple, 21 × 29,7 cm).
    \item \terme{Réseaux sociaux} : dimensions spécifiques (par exemple, 1080 × 1080 px pour Instagram).
\end{itemize}

\subsection{La résolution}

\subsubsection{Définition}

La \terme{résolution} est la quantité de détails qu'une image peut contenir. Elle s'exprime en \terme{points par pouce} (ppp ou dpi, \emph{dots per inch}) pour l'impression, ou en \terme{pixels par pouce} (ppi) pour l'écran.

\subsubsection{Les valeurs courantes}

\begin{itemize}
    \item \terme{72 ppi} : résolution standard pour l'écran (web).
    \item \terme{150 ppi} : résolution moyenne pour l'impression (brouillons, documents).
    \item \terme{300 ppi} : résolution standard pour l'impression professionnelle.
    \item \terme{600 ppi} : résolution haute pour l'impression de qualité (photos, art).
\end{itemize}

\subsubsection{Résolution et qualité}

Plus la résolution est élevée, plus l'image est détaillée, mais plus le fichier est lourd.

\begin{itemize}
    \item \terme{Basse résolution} : image peu détaillée, fichiers légers.
    \item \terme{Haute résolution} : image détaillée, fichiers lourds.
\end{itemize}

\subsection{Relation entre dimensions, résolution et pixels}

\subsubsection{Formule}

La relation entre dimensions physiques, résolution et nombre de pixels est la suivante :

\[
\text{Nombre de pixels} = \text{Dimension physique (pouces)} \times \text{Résolution (ppi)}
\]

\subsubsection{Exemple}

Pour une image de 10 pouces de large à 300 ppi :

\[
10 \times 300 = 3000 \text{ pixels de large}
\]

\subsubsection{Conséquence}

Si on modifie la résolution, on modifie le nombre de pixels (et donc la qualité). Si on modifie les dimensions, on modifie aussi le nombre de pixels.

\subsection{Choisir la bonne résolution}

\subsubsection{Pour l'écran}

\begin{itemize}
    \item \terme{Web} : 72 ppi.
    \item \terme{Écrans haute résolution (Retina)} : 144 ppi ou plus.
    \item \terme{Présentations} : 96 ppi.
\end{itemize}

\subsubsection{Pour l'impression}

\begin{itemize}
    \item \terme{Documents professionnels} : 300 ppi.
    \item \terme{Grands formats} : 150 ppi (affichage à distance).
    \item \terme{Photos haute qualité} : 600 ppi.
\end{itemize}

\subsubsection{Pour les réseaux sociaux}

\begin{itemize}
    \item \terme{Instagram} : 1080 × 1080 px (carré), 1080 × 1350 px (portrait).
    \item \terme{Facebook} : 1200 × 630 px (partage).
    \item \terme{Twitter / X} : 1200 × 675 px.
    \item \terme{LinkedIn} : 1200 × 627 px.
\end{itemize}

\subsection{Les erreurs courantes}

\subsubsection{Travailler en 72 ppi pour l'impression}

La qualité sera insuffisante. Il faut travailler en 300 ppi.

\subsubsection{Agrandir une image de basse résolution}

La pixellisation sera visible. Il faut repartir d'une image haute résolution.

\subsubsection{Confondre dimensions et résolution}

Modifier la résolution change la taille physique (si on garde le nombre de pixels), ou le nombre de pixels (si on garde la taille physique).

\subsubsection{Ne pas anticiper l'usage final}

Il faut déterminer dès le départ l'usage de l'image (écran, impression, réseaux sociaux) pour choisir les bonnes dimensions et la bonne résolution.

\subsection*{Conclusion de la section}

Les \terme{dimensions} d'une image sont sa largeur et sa hauteur, exprimées en pixels ou en unités physiques. La \terme{résolution} est la quantité de détails, exprimée en ppi (écran) ou dpi (impression). La relation entre dimensions, résolution et pixels est donnée par la formule : nombre de pixels = dimension physique × résolution. Pour l'écran, on utilise 72 ppi ; pour l'impression, 300 ppi. Il est essentiel de choisir les bonnes dimensions et la bonne résolution dès le départ, en fonction de l'usage final.

Maintenant que nous maîtrisons les notions de pixels, dimensions et résolution, nous allons, dans la section suivante, étudier les \terme{formats d'image}.

% ============================================================
\section{Formats d'image : JPEG, PNG, TIFF, PSD et autres}
% ============================================================

\subsection{Pourquoi connaître les formats d'image ?}

Connaître les formats d'image est essentiel pour :

\begin{itemize}
    \item \terme{Choisir} le format adapté à l'usage (web, impression, édition).
    \item \terme{Préserver} la qualité de l'image.
    \item \terme{Optimiser} le poids du fichier.
    \item \terme{Éviter} les incompatibilités.
    \item \terme{Préserver} la possibilité de modifier l'image.
\end{itemize}

\subsection{Les principaux formats d'image}

\subsubsection{JPEG (ou JPG)}

\paragraph{Définition} : format d'image matricielle avec \terme{compression avec perte}. Il est très répandu sur le web et en photographie.

\paragraph{Caractéristiques} :
\begin{itemize}
    \item \terme{Compression avec perte} : la qualité est réduite pour diminuer le poids.
    \item \terme{Pas de transparence}.
    \item \terme{Pas de calques}.
    \item \terme{Poids réduit}.
    \item \terme{Compatible} avec tous les navigateurs et logiciels.
\end{itemize}

\paragraph{Usages} :
\begin{itemize}
    \item Photographies pour le web.
    \item Images de partage.
    \item Documents légers.
\end{itemize}

\paragraph{Précautions} :
\begin{itemize}
    \item Éviter de recompter plusieurs fois (perte de qualité cumulative).
    \item Éviter pour les logos, les textes, les aplats de couleur (artefacts).
    \item Ne pas utiliser pour l'impression professionnelle.
\end{itemize}

\subsubsection{PNG}

\paragraph{Définition} : format d'image matricielle avec \terme{compression sans perte}. Il est très utilisé sur le web, notamment pour les images avec transparence.

\paragraph{Caractéristiques} :
\begin{itemize}
    \item \terme{Compression sans perte} : la qualité est préservée.
    \item \terme{Transparence} : supporte la transparence (canal alpha).
    \item \terme{Pas de calques}.
    \item \terme{Poids modéré} (plus lourd que le JPEG).
    \item \terme{Compatible} avec tous les navigateurs.
\end{itemize}

\paragraph{Usages} :
\begin{itemize}
    \item Logos, icônes.
    \item Images avec transparence.
    \item Captures d'écran.
    \item Graphiques avec aplats de couleur.
\end{itemize}

\paragraph{Précautions} :
\begin{itemize}
    \item Éviter pour les photographies (poids plus lourd que le JPEG).
    \item Ne pas utiliser pour l'impression professionnelle (préférer le TIFF).
\end{itemize}

\subsubsection{TIFF}

\paragraph{Définition} : format d'image matricielle avec \terme{compression sans perte}, utilisé pour l'impression professionnelle et l'archivage.

\paragraph{Caractéristiques} :
\begin{itemize}
    \item \terme{Compression sans perte} (ou sans compression).
    \item \terme{Haute qualité} : préserve tous les détails.
    \item \terme{Transparence} : supporte la transparence.
    \item \terme{Calques} : supporte les calques (dans certaines versions).
    \item \terme{Poids très lourd}.
\end{itemize}

\paragraph{Usages} :
\begin{itemize}
    \item Impression professionnelle.
    \item Archivage d'images haute qualité.
    \item Échanges entre professionnels.
\end{itemize}

\paragraph{Précautions} :
\begin{itemize}
    \item Éviter pour le web (poids trop lourd).
    \item Utiliser pour l'impression et l'archivage.
\end{itemize}

\subsubsection{PSD}

\paragraph{Définition} : format natif d'Adobe Photoshop. Il préserve tous les calques, les masques, les réglages.

\paragraph{Caractéristiques} :
\begin{itemize}
    \item \terme{Format natif} de Photoshop.
    \item \terme{Préserve les calques}, les masques, les réglages.
    \item \terme{Modifiable} : on peut rouvrir le fichier et continuer à travailler.
    \item \terme{Poids très lourd}.
    \item \terme{Compatible} uniquement avec Photoshop (et quelques autres logiciels Adobe).
\end{itemize}

\paragraph{Usages} :
\begin{itemize}
    \item Fichiers de travail Photoshop.
    \item Archivage de projets.
    \item Échanges entre designers.
\end{itemize}

\paragraph{Précautions} :
\begin{itemize}
    \item Toujours conserver le PSD comme fichier source.
    \item Ne pas livrer le PSD au client (sauf demande explicite).
    \item Exporter en JPEG, PNG ou TIFF pour la diffusion.
\end{itemize}

\subsubsection{GIF}

\paragraph{Définition} : format d'image matricielle avec \terme{compression sans perte}, limité à 256 couleurs. Il supporte l'animation.

\paragraph{Caractéristiques} :
\begin{itemize}
    \item \terme{Compression sans perte}.
    \item \terme{256 couleurs maximum}.
    \item \terme{Transparence} (1 bit).
    \item \terme{Animation} : supporte les GIF animés.
    \item \terme{Poids léger}.
\end{itemize}

\paragraph{Usages} :
\begin{itemize}
    \item Petites animations (GIF animés).
    \item Icônes simples.
    \item Images avec peu de couleurs.
\end{itemize}

\paragraph{Précautions} :
\begin{itemize}
    \item Éviter pour les photographies (256 couleurs insuffisantes).
    \item Éviter pour les images avec dégradés.
\end{itemize}

\subsubsection{PDF}

\paragraph{Définition} : format de document, qui peut contenir des images matricielles et vectorielles, du texte, des liens.

\paragraph{Caractéristiques} :
\begin{itemize}
    \item \terme{Multi-plateforme} : lisible sur tous les appareils.
    \item \terme{Préserve la mise en page}.
    \item \terme{Peut contenir} du texte, des images, des liens, des formulaires.
    \item \terme{Peut être imprimé} fidèlement.
\end{itemize}

\paragraph{Usages} :
\begin{itemize}
    \item Documents professionnels.
    \item Livraison de fichiers à l'imprimeur.
    \item Présentations.
\end{itemize}

\subsubsection{SVG}

\paragraph{Définition} : format d'image \terme{vectorielle}, basé sur XML. Il est utilisé pour le web.

\paragraph{Caractéristiques} :
\begin{itemize}
    \item \terme{Vectoriel} : pas de perte de qualité à l'agrandissement.
    \item \terme{Léger}.
    \item \terme{Modifiable} avec un éditeur de texte.
    \item \terme{Compatible} avec les navigateurs.
\end{itemize}

\paragraph{Usages} :
\begin{itemize}
    \item Logos, icônes.
    \item Illustrations vectorielles.
    \item Graphiques interactifs.
\end{itemize}

\paragraph{Précautions} :
\begin{itemize}
    \item Éviter pour les photographies.
    \item Utiliser pour les éléments vectoriels.
\end{itemize}

\subsubsection{WebP}

\paragraph{Définition} : format d'image moderne, développé par Google. Il combine les avantages du JPEG et du PNG.

\paragraph{Caractéristiques} :
\begin{itemize}
    \item \terme{Compression avec ou sans perte}.
    \item \terme{Transparence}.
    \item \terme{Animation}.
    \item \terme{Poids réduit}.
\end{itemize}

\paragraph{Usages} :
\begin{itemize}
    \item Images pour le web.
    \item Remplaçant du JPEG et du PNG.
\end{itemize}

\subsection{Tableau récapitulatif des formats}

\begin{table}[h!]
\centering
\begin{tabular}{@{}p{2cm}p{3cm}p{2cm}p{2cm}p{4cm}@{}}
\toprule
\textbf{Format} & \textbf{Type} & \textbf{Compression} & \textbf{Transparence} & \textbf{Usage principal} \\
\midrule
JPEG & Matriciel & Avec perte & Non & Photos web \\
PNG & Matriciel & Sans perte & Oui & Logos, captures \\
TIFF & Matriciel & Sans perte & Oui & Impression, archivage \\
PSD & Matriciel & Sans perte & Oui & Fichiers de travail \\
GIF & Matriciel & Sans perte & Oui & Animations, icônes \\
PDF & Mixte & Variable & Oui & Documents \\
SVG & Vectoriel & Sans perte & Oui & Logos, icônes web \\
WebP & Matriciel & Avec/sans perte & Oui & Images web \\
\bottomrule
\end{tabular}
\caption{Récapitulatif des principaux formats d'image.}
\end{table}

\subsection{Comment choisir un format ?}

\subsubsection{Pour le web}

\begin{itemize}
    \item \terme{Photographies} : JPEG ou WebP.
    \item \terme{Logos, icônes, images avec transparence} : PNG ou SVG.
    \item \terme{Animations} : GIF ou WebP.
\end{itemize}

\subsubsection{Pour l'impression}

\begin{itemize}
    \item \terme{Images haute qualité} : TIFF.
    \item \terme{Documents} : PDF.
    \item \terme{Fichiers de travail} : PSD.
\end{itemize}

\subsubsection{Pour l'archivage}

\begin{itemize}
    \item \terme{Fichiers de travail} : PSD.
    \item \terme{Images haute qualité} : TIFF.
    \item \terme{Documents} : PDF.
\end{itemize}

\subsection*{Conclusion de la section}

Il existe de nombreux formats d'image, chacun avec ses caractéristiques et ses usages. Les principaux sont : JPEG (photos web), PNG (logos, transparence), TIFF (impression, archivage), PSD (fichiers de travail Photoshop), GIF (animations), PDF (documents), SVG (vectoriel web), WebP (web moderne). Le choix d'un format dépend de l'usage : web, impression, archivage. Il est essentiel de connaître ces formats pour travailler efficacement et éviter les erreurs.

Maintenant que nous connaissons les formats d'image, nous allons, dans la section suivante, étudier l'\terme{installation et la configuration du logiciel}.

% ============================================================
\section{Installation et configuration du logiciel}
% ============================================================

\subsection{Les prérequis}

Avant d'installer Photoshop, il faut vérifier que ton ordinateur répond aux prérequis.

\subsubsection{Configuration minimale}

\begin{itemize}
    \item \terme{Système d'exploitation} : Windows 10 (64 bits) ou macOS 10.15 ou supérieur.
    \item \terme{Processeur} : Intel ou AMD multicœur.
    \item \terme{RAM} : 8 Go minimum.
    \item \terme{Espace disque} : 10 Go d'espace libre.
    \item \terme{Écran} : résolution 1280 × 800 minimum.
    \item \terme{Carte graphique} : compatible avec OpenGL 2.0.
\end{itemize}

\subsubsection{Configuration recommandée}

\begin{itemize}
    \item \terme{Système d'exploitation} : Windows 11 ou macOS récent.
    \item \terme{Processeur} : Intel ou AMD récent, multicœur.
    \item \terme{RAM} : 16 Go ou plus.
    \item \terme{Espace disque} : SSD avec 50 Go d'espace libre.
    \item \terme{Écran} : résolution Full HD ou 4K, calibré.
    \item \terme{Carte graphique} : dédiée, récente.
    \item \terme{Tablette graphique} : recommandée pour la retouche et l'illustration.
\end{itemize}

\subsection{L'installation}

\subsubsection{Installation via Creative Cloud}

Photoshop fait partie de la suite Adobe Creative Cloud. Pour l'installer :

\begin{enumerate}
    \item Créer un compte Adobe sur \texttt{adobe.com}.
    \item S'abonner à Creative Cloud (abonnement mensuel ou annuel).
    \item Télécharger et installer l'application \terme{Creative Cloud Desktop}.
    \item Se connecter avec son compte Adobe.
    \item Dans l'application Creative Cloud, sélectionner Photoshop et cliquer sur « Installer ».
    \item Attendre la fin de l'installation.
    \item Lancer Photoshop.
\end{enumerate}

\subsubsection{Autres moyens d'installation}

\begin{itemize}
    \item \terme{Version d'essai} : Adobe propose une version d'essai gratuite de 7 jours.
    \item \terme{Licence éducation} : tarifs réduits pour les étudiants et enseignants.
    \item \terme{Licence entreprise} : pour les entreprises.
\end{itemize}

\subsubsection{Mises à jour}

Photoshop est régulièrement mis à jour. Les mises à jour apportent :

\begin{itemize}
    \item Nouvelles fonctionnalités.
    \item Corrections de bugs.
    \item Améliorations de performance.
    \item Correctifs de sécurité.
\end{itemize}

Il est recommandé de mettre à jour régulièrement, mais avec prudence : parfois, une mise à jour peut introduire des instabilités. Il est conseillé d'attendre quelques jours avant d'installer une mise à jour majeure.

\subsection{La configuration}

\subsubsection{Configuration des préférences}

Une fois Photoshop installé, il est essentiel de configurer les préférences. Pour y accéder : \texttt{Édition > Préférences} (Windows) ou \texttt{Photoshop > Préférences} (macOS).

\paragraph{Général} : langue, historique, interface.

\paragraph{Interface} : thème (clair, sombre), taille de l'interface, couleur de fond.

\paragraph{Performances} : mémoire allouée, accélération GPU, disques de travail.

\paragraph{Curseurs} : apparence des curseurs.

\paragraph{Transparence et gamme de couleurs} : affichage de la transparence, gamme de couleurs.

\paragraph{Unités et règles} : unités de mesure (pixels, centimètres, millimètres, pouces), résolution.

\paragraph{Guides, grille et plans de coupe} : couleur des guides, espacement de la grille.

\paragraph{Plug-ins} : gestion des plug-ins.

\paragraph{Type} : gestion des polices.

\subsubsection{Configuration de la colorimétrie}

Il est essentiel de configurer la colorimétrie pour travailler avec des couleurs fiables.

\paragraph{Espaces colorimétriques} :
\begin{itemize}
    \item \terme{RVB} : sRGB IEC61966-2.1 (standard web).
    \item \terme{CMJN} : selon l'imprimeur (par exemple, ISO Coated v2).
    \item \terme{Niveaux de gris} : Gamma 2.2.
\end{itemize}

\paragraph{Profils de couleur} :
\begin{itemize}
    \item Convertir en profil : selon le support.
    \item Intent : Perceptif (pour les images).
\end{itemize}

\paragraph{Calibration de l'écran} :

Il est essentiel de calibrer son écran avec une sonde (Datacolor Spyder, X-Rite i1Display) pour afficher des couleurs fiables. Sans calibration, les couleurs affichées ne correspondent pas aux couleurs réelles.

\subsubsection{Configuration de l'espace de travail}

Photoshop propose plusieurs espaces de travail prédéfinis :

\begin{itemize}
    \item \terme{Essentiel} : espace de travail par défaut.
    \item \terme{3D} : pour la 3D.
    \item \terme{Motion} : pour l'animation.
    \item \terme{Peinture} : pour la peinture numérique.
    \item \terme{Photographie} : pour la retouche photo.
\end{itemize}

On peut choisir un espace de travail via \texttt{Fenêtre > Espace de travail}. On peut également créer son propre espace de travail et l'enregistrer.

\subsection{Les bonnes pratiques}

\begin{itemize}
    \item \terme{Vérifier les prérequis} avant l'installation.
    \item \terme{Installer} via Creative Cloud.
    \item \terme{Mettre à jour} régulièrement, avec prudence.
    \item \terme{Configurer les préférences} dès le départ.
    \item \terme{Calibrer son écran}.
    \item \terme{Choisir un espace de travail} adapté.
    \item \terme{Créer un espace de travail personnalisé}.
    \item \terme{Sauvegarder ses préférences}.
\end{itemize}

\subsection*{Conclusion de la section}

L'installation et la configuration de Photoshop sont des étapes essentielles. Il faut vérifier les prérequis, installer le logiciel via Creative Cloud, configurer les préférences (interface, performances, unités, etc.), configurer la colorimétrie (espaces, profils, calibration) et choisir un espace de travail adapté. Ces étapes permettent de travailler dans de bonnes conditions et d'éviter les mauvaises surprises.

Nous avons maintenant couvert la première moitié du chapitre 6. Nous avons présenté Photoshop, ses domaines d'utilisation, les images matricielles, les notions de pixels, dimensions et résolution, les formats d'image, l'installation et la configuration. Dans la seconde moitié de ce chapitre, nous aborderons la découverte de l'interface, la barre de menus et les panneaux, les outils essentiels, la création, l'ouverture et l'enregistrement d'un document, les raccourcis clavier fondamentaux, et l'organisation de l'espace de travail.

% ============================================================
\section{Découverte de l'interface}
% ============================================================

\subsection{Pourquoi maîtriser l'interface ?}

L'interface de Photoshop est le \terme{poste de travail} du designer. Elle regroupe tous les outils, tous les panneaux, tous les menus dont tu auras besoin. La maîtriser est indispensable pour travailler efficacement. Un designer qui connaît bien son interface gagne du temps, évite les erreurs, et peut se concentrer sur la création plutôt que sur la recherche des fonctionnalités.

Au début, l'interface de Photoshop peut paraître intimidante : elle comporte de nombreux panneaux, de nombreuses icônes, de nombreux menus. Mais elle est en réalité très logique, une fois qu'on en comprend l'organisation.

\subsection{Les grandes zones de l'interface}

L'interface de Photoshop se compose de plusieurs zones principales.

\subsubsection{La barre de menus}

Située en haut de l'écran, la \terme{barre de menus} contient tous les menus du logiciel : Fichier, Édition, Image, Calque, Texte, Sélection, Filtre, 3D, Affichage, Fenêtre, Aide.

\subsubsection{La barre d'options}

Située juste sous la barre de menus, la \terme{barre d'options} (ou \emph{barre d'outils contextuelle}) affiche les options de l'outil actuellement sélectionné. Elle change selon l'outil.

\subsubsection{La barre d'outils}

Située à gauche de l'écran (par défaut), la \terme{barre d'outils} contient tous les outils de Photoshop : sélection, déplacement, recadrage, pinceau, gomme, texte, etc. Elle peut être affichée en une ou deux colonnes.

\subsubsection{Les panneaux}

Situés à droite de l'écran (par défaut), les \terme{panneaux} regroupent les fonctionnalités par thème : Calques, Couleur, Nuancier, Historique, etc. Ils peuvent être déplacés, regroupés, empilés, réduits en icônes.

\subsubsection{La zone de travail}

Au centre, la \terme{zone de travail} (ou \emph{plan de travail}) affiche le document en cours d'édition. C'est là que l'on crée, retouche, compose.

\subsubsection{La barre d'état}

En bas de la fenêtre, la \terme{barre d'état} affiche des informations sur le document : niveau de zoom, dimensions, poids, etc.

\subsection{Organisation par défaut}

Par défaut, l'interface de Photoshop est organisée ainsi :

\begin{itemize}
    \item \terme{En haut} : barre de menus, barre d'options.
    \item \terme{À gauche} : barre d'outils.
    \item \terme{À droite} : panneaux (Calques, Couleur, etc.).
    \item \terme{Au centre} : zone de travail.
    \item \terme{En bas} : barre d'état.
\end{itemize}

Cette organisation peut être modifiée selon les besoins.

\subsection{Personnalisation de l'interface}

\subsubsection{Déplacer les panneaux}

Les panneaux peuvent être déplacés par \terme{glisser-déposer}. On peut les détacher, les attacher, les empiler, les regrouper.

\subsubsection{Réduire les panneaux}

On peut réduire les panneaux en icônes pour gagner de la place, en cliquant sur la double flèche en haut à droite du panneau.

\subsubsection{Espaces de travail}

Photoshop propose plusieurs espaces de travail prédéfinis (\texttt{Fenêtre > Espace de travail}) :

\begin{itemize}
    \item \terme{Essentiel} : espace par défaut.
    \item \terme{3D} : pour la 3D.
    \item \terme{Motion} : pour l'animation.
    \item \terme{Peinture} : pour la peinture numérique.
    \item \terme{Photographie} : pour la retouche photo.
\end{itemize}

On peut aussi créer son propre espace de travail et l'enregistrer.

\subsubsection{Thème de l'interface}

Photoshop propose deux thèmes : \terme{clair} et \terme{sombre}. On peut les changer dans \texttt{Édition > Préférences > Interface} (Windows) ou \texttt{Photoshop > Préférences > Interface} (macOS). Le thème sombre est souvent préféré pour la retouche photo, car il réduit la fatigue visuelle et évite les interférences avec la perception des couleurs.

\subsection{Les raccourcis d'interface}

\begin{itemize}
    \item \terme{F5} : panneau Nuancier.
    \item \terme{F6} : panneau Couleur.
    \item \terme{F7} : panneau Calques.
    \item \terme{F8} : panneau Informations.
    \item \terme{Tab} : masquer/afficher tous les panneaux.
    \item \terme{Maj + Tab} : masquer/afficher les panneaux sauf la barre d'outils.
\end{itemize}

\subsection{Les bonnes pratiques}

\begin{itemize}
    \item \terme{Choisir un thème} adapté (sombre pour la retouche).
    \item \terme{Organiser ses panneaux} en fonction de son usage.
    \item \terme{Enregistrer son espace de travail} personnalisé.
    \item \terme{Utiliser les raccourcis} pour gagner du temps.
    \item \terme{Ne pas hésiter à réinitialiser} l'espace de travail si nécessaire (\texttt{Fenêtre > Espace de travail > Réinitialiser l'essentiel}).
\end{itemize}

\subsection*{Conclusion de la section}

L'\terme{interface de Photoshop} se compose de plusieurs zones : barre de menus, barre d'options, barre d'outils, panneaux, zone de travail, barre d'état. Elle peut être personnalisée : déplacement des panneaux, réduction en icônes, choix d'un espace de travail, choix du thème. Maîtriser l'interface permet de travailler efficacement et de se concentrer sur la création.

Maintenant que nous avons découvert l'interface, nous allons, dans la section suivante, étudier en détail la \terme{barre de menus et les panneaux}.

% ============================================================
\section{La barre de menus et les panneaux}
% ============================================================

\subsection{La barre de menus}

La barre de menus regroupe tous les menus du logiciel. Chaque menu contient un ensemble de commandes.

\subsubsection{Le menu Fichier}

\begin{itemize}
    \item \terme{Nouveau} : créer un nouveau document.
    \item \terme{Ouvrir} : ouvrir un document existant.
    \item \terme{Ouvrir récent} : ouvrir un document récent.
    \item \terme{Fermer} : fermer le document en cours.
    \item \terme{Enregistrer} : enregistrer le document.
    \item \terme{Enregistrer sous} : enregistrer sous un autre nom ou format.
    \item \terme{Exporter} : exporter le document (JPEG, PNG, etc.).
    \item \terme{Importer} : importer un fichier.
    \item \terme{Automatiser} : automatiser des tâches.
    \item \terme{Imprimer} : imprimer le document.
\end{itemize}

\subsubsection{Le menu Édition}

\begin{itemize}
    \item \terme{Annuler / Rétablir} : annuler ou rétablir une action.
    \item \terme{Copier / Coller} : copier ou coller une sélection.
    \item \terme{Transformations} : rotation, mise à l'échelle, déformation.
    \item \terme{Définir le motif} : créer un motif à partir d'une sélection.
    \item \terme{Remplir} : remplir une sélection avec une couleur, un motif.
    \item \terme{Contour} : ajouter un contour à une sélection.
    \item \terme{Préférences} : configurer le logiciel.
\end{itemize}

\subsubsection{Le menu Image}

\begin{itemize}
    \item \terme{Mode} : changer le mode colorimétrique (RVB, CMJN, etc.).
    \item \terme{Réglages} : luminosité, contraste, courbes, teinte/saturation.
    \item \terme{Taille de l'image} : modifier les dimensions et la résolution.
    \item \terme{Taille de la zone de travail} : modifier la taille du canevas.
    \item \terme{Rotation de l'image} : rotation, symétrie.
    \item \terme{Recadrage} : recadrer l'image.
\end{itemize}

\subsubsection{Le menu Calque}

\begin{itemize}
    \item \terme{Nouveau} : créer un nouveau calque, un groupe, un calque de réglage.
    \item \terme{Dupliquer} : dupliquer un calque.
    \item \terme{Supprimer} : supprimer un calque.
    \item \terme{Fusionner} : fusionner des calques.
    \item \terme{Aplatir l'image} : fusionner tous les calques en un seul.
    \item \terme{Masque} : créer un masque de fusion ou d'écrêtage.
    \item \terme{Calque de réglage} : créer un calque de réglage.
    \item \terme{Style de calque} : ajouter des effets (ombre, contour, etc.).
\end{itemize}

\subsubsection{Le menu Texte}

\begin{itemize}
    \item \terme{Panoramique} : créer un texte.
    \item \terme{Paragraphe} : créer un bloc de texte.
    \item \terme{Vérifier l'orthographe} : corriger l'orthographe.
    \item \terme{Vectoriser} : transformer le texte en forme vectorielle.
\end{itemize}

\subsubsection{Le menu Sélection}

\begin{itemize}
    \item \terme{Tout sélectionner} : sélectionner tout le document.
    \item \terme{Désélectionner} : désélectionner.
    \item \terme{Réactiver la sélection} : réactiver la dernière sélection.
    \item \terme{Inverser} : inverser la sélection.
    \item \terme{Modifier} : modifier la sélection (contour, lissage, etc.).
    \item \terme{Plage de couleurs} : sélectionner par plage de couleurs.
\end{itemize}

\subsubsection{Le menu Filtre}

\begin{itemize}
    \item \terme{Flou} : flou gaussien, flou de mouvement, etc.
    \item \terme{Netteté} : accentuer, renforcer la netteté.
    \item \terme{Distorsion} : déformation, ondulation, etc.
    \item \terme{Bruit} : ajout ou réduction de bruit.
    \item \terme{Esthétique} : effets artistiques.
    \item \terme{Rendu} : nuages, éclairage, etc.
    \item \terme{Galerie de filtres} : appliquer plusieurs filtres.
    \item \terme{Intelligence artificielle} : filtres IA (Neural Filters).
\end{itemize}

\subsubsection{Le menu 3D}

\begin{itemize}
    \item Création et édition d'objets 3D.
    \item (Ce menu n'existe que dans les versions compatibles 3D.)
\end{itemize}

\subsubsection{Le menu Affichage}

\begin{itemize}
    \item \terme{Zoom} : zoom avant, zoom arrière, ajuster à l'écran.
    \item \terme{Repères} : règles, guides, grille.
    \item \terme{Aligner sur} : aligner sur les guides, la grille.
    \item \terme{Afficher} : afficher ou masquer des éléments.
\end{itemize}

\subsubsection{Le menu Fenêtre}

\begin{itemize}
    \item \terme{Espaces de travail} : choisir un espace de travail.
    \item \terme{Panneaux} : afficher ou masquer des panneaux.
    \item \terme{Nouvelle fenêtre} : ouvrir une nouvelle fenêtre.
    \item \terme{Organiser} : organiser les fenêtres.
\end{itemize}

\subsubsection{Le menu Aide}

\begin{itemize}
    \item \terme{Photoshop Help} : aide en ligne.
    \item \terme{À propos de Photoshop} : informations sur la version.
\end{itemize}

\subsection{Les panneaux}

Les panneaux sont des fenêtres flottantes ou ancrées qui regroupent des fonctionnalités par thème. Voici les principaux.

\subsubsection{Le panneau Calques}

Le panneau \terme{Calques} (F7) est le panneau le plus important de Photoshop. Il affiche la liste des calques du document, et permet de :

\begin{itemize}
    \item Créer, dupliquer, supprimer des calques.
    \item Renommer, organiser, regrouper des calques.
    \item Afficher ou masquer des calques.
    \item Verrouiller des calques.
    \item Changer l'opacité, les modes de fusion.
    \item Ajouter des masques, des styles de calque, des calques de réglage.
\end{itemize}

\subsubsection{Le panneau Couleur}

Le panneau \terme{Couleur} (F6) permet de choisir et de mélanger des couleurs. Il propose plusieurs modes : TSL, RVB, CMJN, etc.

\subsubsection{Le panneau Nuancier}

Le panneau \terme{Nuancier} (F5) regroupe des couleurs prédéfinies ou enregistrées. On peut y ajouter ses propres couleurs.

\subsubsection{Le panneau Historique}

Le panneau \terme{Historique} permet de revenir en arrière dans les actions effectuées. On peut annuler jusqu'à 50 étapes (par défaut). On peut aussi créer des instantanés pour sauvegarder un état.

\subsubsection{Le panneau Propriétés}

Le panneau \terme{Propriétés} affiche les propriétés du calque ou de l'outil sélectionné. Il permet d'ajuster rapidement les paramètres.

\subsubsection{Le panneau Réglages}

Le panneau \terme{Réglages} (ou \emph{Ajustements}) permet d'ajouter des calques de réglage : luminosité/contraste, niveaux, courbes, teinte/saturation, etc.

\subsubsection{Le panneau Caractère et Paragraphe}

Ces panneaux permettent de gérer la typographie : police, taille, graisse, style, interlignage, crénage, alignement, etc.

\subsubsection{Le panneau Chemin}

Le panneau \terme{Chemin} affiche les tracés vectoriels créés avec l'outil Plume. Il permet de les enregistrer, de les convertir en sélection, etc.

\subsubsection{Le panneau Formes}

Le panneau \terme{Formes} (Bibliothèque de formes) propose des formes prédéfinies : flèches, étoiles, bulles, etc.

\subsubsection{Le panneau Bibliothèques}

Les \terme{Bibliothèques} Creative Cloud permettent de synchroniser des ressources (couleurs, formes, styles) entre les appareils et les logiciels Adobe.

\subsection{Organiser ses panneaux}

\subsubsection{Regrouper}

On peut regrouper plusieurs panneaux dans une même colonne, en les glissant l'un sur l'autre.

\subsubsection{Empiler}

On peut empiler des panneaux pour gagner de la place.

\subsubsection{Détacher}

On peut détacher un panneau pour le rendre flottant.

\subsubsection{Enregistrer un espace de travail}

On peut enregistrer son organisation de panneaux dans un espace de travail personnalisé (\texttt{Fenêtre > Espace de travail > Nouvel espace de travail}).

\subsection{Les bonnes pratiques}

\begin{itemize}
    \item \terme{Garder le panneau Calques} toujours visible.
    \item \terme{Garder les panneaux Couleur et Nuancier} à portée de main.
    \item \terme{Utiliser le panneau Historique} pour revenir en arrière.
    \item \terme{Regrouper les panneaux} par usage.
    \item \terme{Enregistrer son espace de travail} personnalisé.
    \item \terme{Utiliser les raccourcis} (F5, F6, F7, F8).
\end{itemize}

\subsection*{Conclusion de la section}

La \terme{barre de menus} regroupe tous les menus du logiciel : Fichier, Édition, Image, Calque, Texte, Sélection, Filtre, 3D, Affichage, Fenêtre, Aide. Les \terme{panneaux} regroupent les fonctionnalités par thème : Calques, Couleur, Nuancier, Historique, Propriétés, Réglages, Caractère, etc. Maîtriser la barre de menus et les panneaux est essentiel pour travailler efficacement. On peut organiser ses panneaux, les regrouper, les empiler, les détacher, et enregistrer son espace de travail.

Maintenant que nous connaissons la barre de menus et les panneaux, nous allons, dans la section suivante, étudier les \terme{outils essentiels} de Photoshop.

% ============================================================
\section{Les outils essentiels}
% ============================================================

\subsection{La barre d'outils}

La barre d'outils regroupe tous les outils de Photoshop. Elle est organisée par catégories. Certains outils sont masqués derrière d'autres : on y accède en cliquant longuement sur l'icône, ou en utilisant le raccourci clavier correspondant.

\subsection{Les outils de sélection}

\subsubsection{L'outil Déplacement (V)}

L'outil \terme{Déplacement} permet de déplacer des calques, des sélections, des objets.

\subsubsection{L'outil Rectangle de sélection (M)}

L'outil \terme{Rectangle de sélection} permet de sélectionner une zone rectangulaire. En maintenant \texttt{Maj}, on obtient un carré.

\subsubsection{L'outil Ellipse de sélection (M)}

L'outil \terme{Ellipse de sélection} permet de sélectionner une zone elliptique. En maintenant \texttt{Maj}, on obtient un cercle.

\subsubsection{L'outil Lasso (L)}

L'outil \terme{Lasso} permet de dessiner une sélection à main levée.

\subsubsection{L'outil Lasso polygonal (L)}

L'outil \terme{Lasso polygonal} permet de dessiner une sélection avec des segments droits.

\subsubsection{L'outil Lasso magnétique (L)}

L'outil \terme{Lasso magnétique} détecte automatiquement les contours des objets.

\subsubsection{L'outil Sélection rapide (W)}

L'outil \terme{Sélection rapide} sélectionne automatiquement des zones en fonction de la couleur et de la texture.

\subsubsection{L'outil Baguette magique (W)}

L'outil \terme{Baguette magique} sélectionne des zones de couleur similaire en un clic.

\subsubsection{L'outil Sélection d'objet (W)}

L'outil \terme{Sélection d'objet} (ou \emph{Select Subject}) sélectionne automatiquement le sujet principal d'une image grâce à l'IA.

\subsection{Les outils de recadrage et de découpe}

\subsubsection{L'outil Recadrage (C)}

L'outil \terme{Recadrage} permet de recadrer une image. On peut choisir un ratio (carré, 16:9, etc.) ou recadrer librement.

\subsubsection{L'outil Perspective (C)}

L'outil \terme{Perspective} permet de corriger la perspective d'une image (redresser un bâtiment, par exemple).

\subsubsection{L'outil Découpe (C)}

L'outil \terme{Découpe} permet de découper une image en plusieurs parties.

\subsubsection{L'outil Sélection de segment (C)}

Utilisé pour le web, il permet de découper une image en zones cliquables.

\subsection{Les outils de retouche}

\subsubsection{L'outil Correcteur (J)}

L'outil \terme{Correcteur} (Spot Healing Brush) permet de supprimer les imperfections en un clic, en les remplaçant par les pixels environnants.

\subsubsection{L'outil Correcteur par plage (J)}

L'outil \terme{Correcteur par plage} permet de supprimer des imperfections en utilisant une zone source définie.

\subsubsection{L'outil Pièce (J)}

L'outil \terme{Pièce} (Patch) permet de remplacer une zone par une autre, en respectant la texture et la couleur.

\subsubsection{L'outil Œil rouge (J)}

L'outil \terme{Œil rouge} permet de corriger les yeux rouges.

\subsubsection{L'outil Tampon (S)}

L'outil \terme{Tampon} (Clone Stamp) permet de dupliquer une partie d'une image sur une autre.

\subsubsection{L'outil Motif (S)}

L'outil \terme{Motif} (Pattern Stamp) permet de peindre avec un motif.

\subsubsection{L'outil Gomme (E)}

L'outil \terme{Gomme} permet d'effacer des pixels.

\subsubsection{L'outil Gomme magique (E)}

L'outil \terme{Gomme magique} efface les pixels de couleur similaire.

\subsubsection{L'outil Gomme d'arrière-plan (E)}

L'outil \terme{Gomme d'arrière-plan} efface l'arrière-plan en préservant le sujet.

\subsubsection{L'outil Netteté (R)}

L'outil \terme{Netteté} (Sharpen) permet d'accentuer la netteté localement.

\subsubsection{L'outil Goutte d'eau (R)}

L'outil \terme{Goutte d'eau} (Smudge) permet d'étirer les pixels pour créer un effet de flou.

\subsubsection{L'outil Densité (R)}

L'outil \terme{Densité} (Dodge) permet d'éclaircir ou d'assombrir localement (Dodge / Burn).

\subsection{Les outils de peinture}

\subsubsection{L'outil Pinceau (B)}

L'outil \terme{Pinceau} permet de peindre avec une couleur. On peut choisir la taille, la dureté, l'opacité, le débit, la forme.

\subsubsection{L'outil Crayon (B)}

L'outil \terme{Crayon} peint avec des bords nets (pas d'anti-aliasing).

\subsubsection{L'outil Remplacement de couleur (B)}

L'outil \terme{Remplacement de couleur} peint avec une couleur qui remplace la couleur existante.

\subsubsection{L'outil Mélangeur (B)}

L'outil \terme{Mélangeur} (Mixer Brush) simule la peinture à l'huile.

\subsubsection{L'outil Pot de peinture (G)}

L'outil \terme{Pot de peinture} remplit une zone avec une couleur.

\subsubsection{L'outil Dégradé (G)}

L'outil \terme{Dégradé} crée un dégradé linéaire, radial, angulaire, réfléchi ou en losange.

\subsection{Les outils de dessin vectoriel}

\subsubsection{L'outil Plume (P)}

L'outil \terme{Plume} permet de créer des tracés vectoriels (courbes de Bézier).

\subsubsection{L'outil Plume de forme libre (P)}

L'outil \terme{Plume de forme libre} permet de dessiner à main levée.

\subsubsection{L'outil Ajouter un point d'ancrage (P)}

Permet d'ajouter un point d'ancrage sur un tracé existant.

\subsubsection{L'outil Supprimer un point d'ancrage (P)}

Permet de supprimer un point d'ancrage.

\subsubsection{L'outil Conversion de point (P)}

Permet de convertir un point d'angle en point lisse, ou inversement.

\subsubsection{L'outil Rectangle (U)}

Permet de créer des rectangles vectoriels.

\subsubsection{L'outil Rectangle arrondi (U)}

Permet de créer des rectangles à coins arrondis.

\subsubsection{L'outil Ellipse (U)}

Permet de créer des ellipses vectorielles.

\subsubsection{L'outil Polygone (U)}

Permet de créer des polygones vectoriels.

\subsubsection{L'outil Trait (U)}

Permet de créer des traits vectoriels.

\subsubsection{L'outil Forme personnalisée (U)}

Permet de créer des formes prédéfinies (flèches, étoiles, etc.).

\subsection{Les outils de texte}

\subsubsection{L'outil Texte horizontal (T)}

Permet de créer un texte horizontal.

\subsubsection{L'outil Texte vertical (T)}

Permet de créer un texte vertical.

\subsubsection{L'outil Texte masque horizontal (T)}

Permet de créer une sélection en forme de texte.

\subsubsection{L'outil Texte masque vertical (T)}

Permet de créer une sélection en forme de texte vertical.

\subsection{Les outils de navigation}

\subsubsection{L'outil Main (H)}

Permet de déplacer la vue dans la zone de travail.

\subsubsection{L'outil Zoom (Z)}

Permet de zoomer avant ou arrière.

\subsubsection{L'outil Rotation de la vue (R)}

Permet de faire pivoter la vue (sans modifier l'image).

\subsection{Les outils de couleur}

\subsubsection{La couleur de premier plan}

La \terme{couleur de premier plan} est la couleur utilisée par défaut pour peindre, remplir, etc.

\subsubsection{La couleur d'arrière-plan}

La \terme{couleur d'arrière-plan} est la couleur utilisée pour les dégradés, les gommes, etc.

\subsubsection{Le sélecteur de couleur}

Cliquer sur la couleur de premier plan ou d'arrière-plan ouvre le \terme{sélecteur de couleur}, qui permet de choisir une couleur précise.

\subsubsection{L'icône de permutation}

L'icône en forme de flèche courbe permet de permuter les couleurs de premier plan et d'arrière-plan (raccourci \texttt{X}).

\subsubsection{L'icône Noir/Blanc}

L'icône en forme de carré noir et blanc permet de revenir aux couleurs par défaut (raccourci \texttt{D}).

\subsection{Les raccourcis des outils}

\begin{table}[h!]
\centering
\begin{tabular}{@{}ll@{}}
\toprule
\textbf{Outil} & \textbf{Raccourci} \\
\midrule
Déplacement & V \\
Rectangle de sélection & M \\
Lasso & L \\
Sélection rapide & W \\
Recadrage & C \\
Correcteur & J \\
Tampon & S \\
Gomme & E \\
Pinceau & B \\
Pot de peinture / Dégradé & G \\
Plume & P \\
Forme & U \\
Texte & T \\
Main & H \\
Zoom & Z \\
Couleur par défaut & D \\
Permuter les couleurs & X \\
\bottomrule
\end{tabular}
\caption{Raccourcis clavier des outils essentiels.}
\end{table}

\subsection{Les bonnes pratiques}

\begin{itemize}
    \item \terme{Apprendre les raccourcis} progressivement.
    \item \terme{Utiliser l'outil adapté} à la tâche.
    \item \terme{Ne pas hésiter à expérimenter}.
    \item \terme{Combiner les outils} pour obtenir le résultat souhaité.
    \item \terme{Travailler de manière non destructive} (calques, masques).
\end{itemize}

\subsection*{Conclusion de la section}

Photoshop propose de nombreux outils, classés par catégories : sélection, recadrage, retouche, peinture, dessin vectoriel, texte, navigation, couleur. Chaque outil a une fonction spécifique et un raccourci clavier. Maîtriser les outils essentiels est indispensable pour travailler efficacement. Les raccourcis clavier permettent de gagner du temps et de se concentrer sur la création.

Maintenant que nous connaissons les outils essentiels, nous allons, dans la section suivante, étudier la \terme{création, l'ouverture et l'enregistrement d'un document}.

% ============================================================
\section{Création, ouverture et enregistrement d'un document}
% ============================================================

\subsection{Créer un nouveau document}

\subsubsection{La commande Nouveau}

Pour créer un nouveau document : \texttt{Fichier > Nouveau} (raccourci \texttt{Ctrl + N} sur Windows, \texttt{Cmd + N} sur macOS).

\subsubsection{La boîte de dialogue}

La boîte de dialogue \emph{Nouveau document} propose plusieurs paramètres.

\paragraph{Nom} : donner un nom au document.

\paragraph{Format prédéfini} : choisir un format prédéfini (A4, Web, Mobile, etc.).

\paragraph{Largeur et hauteur} : définir les dimensions du document.

\paragraph{Unités} : pixels, centimètres, millimètres, pouces, points, picas.

\paragraph{Résolution} : définir la résolution (72 ppi pour l'écran, 300 ppi pour l'impression).

\paragraph{Orientation} : portrait ou paysage.

\paragraph{Couleur de fond} : blanc, noir, arrière-plan, transparent.

\paragraph{Profil colorimétrique} : RVB, CMJN, niveaux de gris.

\paragraph{Profondeur de bits} : 8 bits, 16 bits, 32 bits.

\paragraph{Profil colorimétrique} : sRGB, Adobe RGB, etc.

\subsubsection{Les formats prédéfinis}

Photoshop propose de nombreux formats prédéfinis :

\begin{itemize}
    \item \terme{Photo} : tailles courantes pour la photo.
    \item \terme{Web} : 1920 × 1080, 1280 × 720, etc.
    \item \terme{Impression} : A4, A3, A5, etc.
    \item \terme{Mobile} : iPhone, iPad, etc.
    \item \terme{Film et vidéo} : 4K, Full HD, etc.
\end{itemize}

\subsubsection{Les bonnes pratiques}

\begin{itemize}
    \item \terme{Définir les bonnes dimensions} dès le départ.
    \item \terme{Choisir la bonne résolution} (72 ppi pour l'écran, 300 ppi pour l'impression).
    \item \terme{Choisir le bon profil colorimétrique} (sRGB pour le web, Adobe RGB pour la photo, CMJN pour l'impression).
    \item \terme{Travailler en 8 bits} pour la plupart des projets (16 bits pour la retouche avancée).
    \item \terme{Nommer le document} de manière claire.
\end{itemize}

\subsection{Ouvrir un document}

\subsubsection{La commande Ouvrir}

Pour ouvrir un document : \texttt{Fichier > Ouvrir} (raccourci \texttt{Ctrl + O} sur Windows, \texttt{Cmd + O} sur macOS).

\subsubsection{Les formats supportés}

Photoshop supporte de nombreux formats : PSD, JPEG, PNG, TIFF, GIF, PDF, etc.

\subsubsection{L'ouverture récente}

Pour ouvrir un document récent : \texttt{Fichier > Ouvrir récent}.

\subsubsection{L'ouverture en tant qu'objet dynamique}

On peut ouvrir un fichier en tant qu'objet dynamique (\texttt{Fichier > Ouvrir en tant qu'objet dynamique}), ce qui permet de le modifier sans perte.

\subsubsection{L'importation}

Pour importer un fichier dans un document existant : \texttt{Fichier > Importer} ou \texttt{Fichier > Importer et placer}.

\subsection{Enregistrer un document}

\subsubsection{La commande Enregistrer}

Pour enregistrer un document : \texttt{Fichier > Enregistrer} (raccourci \texttt{Ctrl + S} sur Windows, \texttt{Cmd + S} sur macOS).

Si le document est nouveau, la boîte de dialogue \emph{Enregistrer sous} s'ouvre.

\subsubsection{La commande Enregistrer sous}

Pour enregistrer sous un autre nom ou format : \texttt{Fichier > Enregistrer sous} (raccourci \texttt{Ctrl + Maj + S} sur Windows, \texttt{Cmd + Maj + S} sur macOS).

\subsubsection{Les formats d'enregistrement}

\begin{itemize}
    \item \terme{PSD} : format natif de Photoshop, préserve les calques.
    \item \terme{PSB} : format pour les très grands documents.
    \item \terme{TIFF} : format pour l'impression.
    \item \terme{JPEG} : format pour le web (avec perte).
    \item \terme{PNG} : format pour le web (avec transparence).
    \item \terme{PDF} : format de document.
    \item \terme{GIF} : format pour les animations.
\end{itemize}

\subsubsection{La commande Exporter}

Pour exporter un document : \texttt{Fichier > Exporter > Exporter sous} ou \texttt{Fichier > Exporter > Exporter rapidement en tant que}.

\paragraph{Exporter sous} : permet de choisir le format, la qualité, les dimensions.

\paragraph{Exporter rapidement en tant que} : permet d'exporter dans un format prédéfini (PNG, JPEG, GIF).

\subsubsection{Le format pour le web}

Pour optimiser une image pour le web : \texttt{Fichier > Exporter > Enregistrer pour le web (hérité)}. Cette fonction permet de choisir le format, la qualité, les dimensions, et de visualiser le résultat.

\subsection{Les bonnes pratiques}

\begin{itemize}
    \item \terme{Toujours conserver le PSD} comme fichier source.
    \item \terme{Enregistrer régulièrement} pour éviter les pertes.
    \item \terme{Utiliser la commande Enregistrer sous} pour créer des versions.
    \item \terme{Exporter dans le bon format} selon l'usage (web, impression).
    \item \terme{Optimiser les images} pour le web (réduire le poids).
    \item \terme{Organiser ses fichiers} dans des dossiers clairs.
\end{itemize}

\subsection*{Conclusion de la section}

Créer, ouvrir et enregistrer un document sont des opérations fondamentales dans Photoshop. La création passe par la boîte de dialogue \emph{Nouveau}, qui permet de définir les dimensions, la résolution, le profil colorimétrique. L'ouverture se fait par la commande \emph{Ouvrir} ou \emph{Importer}. L'enregistrement se fait par les commandes \emph{Enregistrer}, \emph{Enregistrer sous}, \emph{Exporter}. Il est essentiel de connaître les formats et de choisir le bon selon l'usage.

Maintenant que nous savons créer, ouvrir et enregistrer un document, nous allons, dans la section suivante, étudier les \terme{raccourcis clavier fondamentaux}.

% ============================================================
\section{Les raccourcis clavier fondamentaux}
% ============================================================

\subsection{Pourquoi apprendre les raccourcis ?}

Les raccourcis clavier sont essentiels pour travailler efficacement dans Photoshop. Ils permettent de :

\begin{itemize}
    \item \terme{Gagner du temps} : éviter de chercher dans les menus.
    \item \terme{Rester concentré} : ne pas interrompre le travail.
    \item \terme{Améliorer la fluidité} : enchaîner les actions rapidement.
    \item \terme{Travailler comme un professionnel} : les pros utilisent tous les raccourcis.
\end{itemize}

Au début, il peut être difficile de les mémoriser. Mais avec la pratique, ils deviennent automatiques.

\subsection{Les raccourcis essentiels}

\subsubsection{Fichier}

\begin{table}[h!]
\centering
\begin{tabular}{@{}ll@{}}
\toprule
\textbf{Action} & \textbf{Raccourci} \\
\midrule
Nouveau document & Ctrl + N (Cmd + N) \\
Ouvrir & Ctrl + O (Cmd + O) \\
Enregistrer & Ctrl + S (Cmd + S) \\
Enregistrer sous & Ctrl + Maj + S (Cmd + Maj + S) \\
Exporter sous & Ctrl + Alt + Maj + S (Cmd + Opt + Maj + S) \\
Fermer & Ctrl + W (Cmd + W) \\
Imprimer & Ctrl + P (Cmd + P) \\
\bottomrule
\end{tabular}
\caption{Raccourcis de gestion de fichiers.}
\end{table}

\subsubsection{Édition}

\begin{table}[h!]
\centering
\begin{tabular}{@{}ll@{}}
\toprule
\textbf{Action} & \textbf{Raccourci} \\
\midrule
Annuler & Ctrl + Z (Cmd + Z) \\
Rétablir & Ctrl + Maj + Z (Cmd + Maj + Z) \\
Copier & Ctrl + C (Cmd + C) \\
Coller & Ctrl + V (Cmd + V) \\
Couper & Ctrl + X (Cmd + X) \\
Tout sélectionner & Ctrl + A (Cmd + A) \\
Désélectionner & Ctrl + D (Cmd + D) \\
Réactiver la sélection & Ctrl + Maj + D (Cmd + Maj + D) \\
Inverser la sélection & Ctrl + Maj + I (Cmd + Maj + I) \\
Transformation libre & Ctrl + T (Cmd + T) \\
Remplir & Maj + F5 \\
\bottomrule
\end{tabular}
\caption{Raccourcis d'édition.}
\end{table}

\subsubsection{Outils}

\begin{table}[h!]
\centering
\begin{tabular}{@{}ll@{}}
\toprule
\textbf{Outil} & \textbf{Raccourci} \\
\midrule
Déplacement & V \\
Sélection rectangle & M \\
Lasso & L \\
Sélection rapide & W \\
Recadrage & C \\
Correcteur & J \\
Tampon & S \\
Gomme & E \\
Pinceau & B \\
Pot de peinture / Dégradé & G \\
Plume & P \\
Forme & U \\
Texte & T \\
Main & H \\
Zoom & Z \\
Couleur par défaut & D \\
Permuter les couleurs & X \\
\bottomrule
\end{tabular}
\caption{Raccourcis des outils.}
\end{table}

\subsubsection{Calques}

\begin{table}[h!]
\centering
\begin{tabular}{@{}ll@{}}
\toprule
\textbf{Action} & \textbf{Raccourci} \\
\midrule
Nouveau calque & Ctrl + Maj + N (Cmd + Maj + N) \\
Dupliquer le calque & Ctrl + J (Cmd + J) \\
Fusionner vers le bas & Ctrl + E (Cmd + E) \\
Aplatir l'image & Ctrl + Maj + E (Cmd + Maj + E) \\
Groupe de calques & Ctrl + G (Cmd + G) \\
Créer un masque & Bouton en bas du panneau Calques \\
Calque de réglage & Bouton en bas du panneau Calques \\
\bottomrule
\end{tabular}
\caption{Raccourcis de gestion des calques.}
\end{table}

\subsubsection{Affichage}

\begin{table}[h!]
\centering
\begin{tabular}{@{}ll@{}}
\toprule
\textbf{Action} & \textbf{Raccourci} \\
\midrule
Zoom avant & Ctrl + + \\
Zoom arrière & Ctrl + - \\
Ajuster à l'écran & Ctrl + 0 (zéro) \\
Taille réelle & Ctrl + 1 \\
Masquer les panneaux & Tab \\
Masquer sauf la barre d'outils & Maj + Tab \\
Afficher/masquer les règles & Ctrl + R \\
Afficher/masquer les guides & Ctrl + ; \\
Afficher/masquer la grille & Ctrl + ' \\
\bottomrule
\end{tabular}
\caption{Raccourcis d'affichage.}
\end{table}

\subsubsection{Sélection}

\begin{table}[h!]
\centering
\begin{tabular}{@{}ll@{}}
\toprule
\textbf{Action} & \textbf{Raccourci} \\
\midrule
Ajouter à la sélection & Maj + clic \\
Soustraire de la sélection & Alt + clic (Opt + clic) \\
Intersection & Maj + Alt + clic \\
Déplacer la sélection & Flèches \\
Copier la sélection & Ctrl + C \\
Couper la sélection & Ctrl + X \\
Contour progressif & Maj + F6 \\
\bottomrule
\end{tabular}
\caption{Raccourcis de sélection.}
\end{table}

\subsubsection{Pinceau}

\begin{table}[h!]
\centering
\begin{tabular}{@{}ll@{}}
\toprule
\textbf{Action} & \textbf{Raccourci} \\
\midrule
Diminuer la taille & [ \\
Augmenter la taille & ] \\
Diminuer la dureté & Maj + [ \\
Augmenter la dureté & Maj + ] \\
Opacité 100\% & 0 \\
Opacité 10\% & 1 \\
\bottomrule
\end{tabular}
\caption{Raccourcis du pinceau.}
\end{table}

\subsection{Les bonnes pratiques}

\begin{itemize}
    \item \terme{Apprendre progressivement} : ne pas essayer de tout mémoriser d'un coup.
    \item \terme{Pratiquer régulièrement} : utiliser les raccourcis dès que possible.
    \item \terme{Créer ses propres raccourcis} : Photoshop permet de personnaliser les raccourcis (\texttt{Édition > Raccourcis clavier}).
    \item \terme{Afficher les raccourcis} : on peut afficher les raccourcis dans les menus.
    \item \terme{Ne pas hésiter à utiliser les menus} au début.
\end{itemize}

\subsection*{Conclusion de la section}

Les \terme{raccourcis clavier} sont essentiels pour travailler efficacement dans Photoshop. Ils couvrent la gestion de fichiers, l'édition, les outils, les calques, l'affichage, la sélection, le pinceau. Il faut les apprendre progressivement et les pratiquer régulièrement. Photoshop permet également de personnaliser ses raccourcis.

Maintenant que nous connaissons les raccourcis fondamentaux, nous allons, dans la dernière section de ce chapitre, étudier l'\terme{organisation de l'espace de travail}.

% ============================================================
\section{Organisation de l'espace de travail}
% ============================================================

\subsection{Pourquoi organiser son espace de travail ?}

Un espace de travail bien organisé permet de :

\begin{itemize}
    \item \terme{Gagner du temps} : avoir les outils et panneaux à portée de main.
    \item \terme{Travailler efficacement} : éviter les manipulations inutiles.
    \item \terme{Réduire la fatigue} : avoir un environnement agréable.
    \item \terme{Se concentrer} : éviter les distractions.
    \item \terme{Adapter} l'espace à chaque projet.
\end{itemize}

\subsection{Les espaces de travail prédéfinis}

Photoshop propose plusieurs espaces de travail prédéfinis (\texttt{Fenêtre > Espace de travail}) :

\begin{itemize}
    \item \terme{Essentiel} : espace par défaut, polyvalent.
    \item \terme{3D} : pour la modélisation 3D.
    \item \terme{Graphisme et web} : pour le design web.
    \item \terme{Motion} : pour l'animation.
    \item \terme{Peinture} : pour la peinture numérique.
    \item \terme{Photographie} : pour la retouche photo.
\end{itemize}

On peut choisir un espace de travail selon le projet.

\subsection{Créer un espace de travail personnalisé}

\subsubsection{Étape 1 : Organiser les panneaux}

Déplacer, regrouper, empiler les panneaux selon ses besoins.

\subsubsection{Étape 2 : Configurer la barre d'outils}

Afficher les outils que l'on utilise fréquemment, masquer les autres.

\subsubsection{Étape 3 : Enregistrer l'espace de travail}

\texttt{Fenêtre > Espace de travail > Nouvel espace de travail}. Donner un nom à l'espace de travail.

\subsubsection{Étape 4 : Réinitialiser si nécessaire}

\texttt{Fenêtre > Espace de travail > Réinitialiser <nom de l'espace>}.

\subsection{Les bonnes pratiques d'organisation}

\subsubsection{Garder le panneau Calques visible}

Le panneau Calques est le plus utilisé : il doit toujours être visible.

\subsubsection{Garder les panneaux Couleur et Nuancier accessibles}

Ces panneaux sont utiles pour choisir et gérer les couleurs.

\subsubsection{Utiliser le panneau Historique}

Le panneau Historique permet de revenir en arrière : il doit être accessible.

\subsubsection{Regrouper les panneaux par usage}

Par exemple : Calques + Historique + Propriétés dans une colonne ; Couleur + Nuancier + Réglages dans une autre.

\subsubsection{Utiliser le mode plein écran}

Le mode plein écran (\texttt{F}) permet de masquer les panneaux et de se concentrer sur l'image. On peut basculer entre les modes avec \texttt{F}.

\subsubsection{Personnaliser la barre d'outils}

On peut personnaliser la barre d'outils en cliquant sur les trois points en bas de la barre (\texttt{Modifier la barre d'outils}). On peut y ajouter ou supprimer des outils.

\subsection{Les raccourcis d'organisation}

\begin{itemize}
    \item \terme{Tab} : masquer/afficher tous les panneaux.
    \item \terme{Maj + Tab} : masquer/afficher les panneaux sauf la barre d'outils.
    \item \terme{F} : basculer entre les modes d'affichage (standard, plein écran avec menus, plein écran).
    \item \terme{Ctrl + 0} : ajuster l'image à l'écran.
    \item \terme{Ctrl + 1} : afficher à taille réelle.
\end{itemize}

\subsection{Les bonnes habitudes}

\subsubsection{Commencer chaque projet par organiser son espace}

Avant de commencer un projet, prendre quelques secondes pour organiser son espace de travail en fonction du projet.

\subsubsection{Enregistrer plusieurs espaces de travail}

On peut enregistrer plusieurs espaces de travail : un pour la retouche photo, un pour la création publicitaire, un pour l'illustration, etc.

\subsubsection{Ne pas hésiter à réinitialiser}

Si l'espace de travail est perturbé, on peut le réinitialiser.

\subsubsection{Mettre à jour ses espaces}

Au fur et à mesure que l'on évolue, on peut mettre à jour ses espaces de travail.

\subsection*{Conclusion de la section}

L'\terme{organisation de l'espace de travail} est essentielle pour travailler efficacement dans Photoshop. On peut utiliser les espaces de travail prédéfinis, ou créer son propre espace personnalisé. Il est recommandé de garder le panneau Calques visible, de regrouper les panneaux par usage, d'utiliser le mode plein écran, de personnaliser la barre d'outils. Les bonnes habitudes consistent à organiser son espace avant chaque projet, à enregistrer plusieurs espaces, à les réinitialiser si nécessaire.

% ============================================================
\section*{Conclusion du chapitre}
\addcontentsline{toc}{section}{Conclusion du chapitre}
% ============================================================

Nous voici au terme du chapitre 6. Nous avons découvert Adobe Photoshop sous tous ses aspects :

\begin{itemize}
    \item \terme{Présentation} : logiciel de création et de retouche d'images matricielles, standard professionnel.
    \item \terme{Domaines d'utilisation} : retouche photo, photomontage, publicité, web, illustration, mockups, restauration.
    \item \terme{Images matricielles} : composées de pixels, précision mais perte de qualité à l'agrandissement.
    \item \terme{Pixels, dimensions, résolution} : notions fondamentales.
    \item \terme{Formats d'image} : JPEG, PNG, TIFF, PSD, GIF, PDF, SVG, WebP.
    \item \terme{Installation et configuration} : prérequis, installation, préférences, colorimétrie.
    \item \terme{Interface} : barre de menus, barre d'options, barre d'outils, panneaux, zone de travail.
    \item \terme{Barre de menus et panneaux} : Fichier, Édition, Image, Calque, etc. ; Calques, Couleur, Historique, etc.
    \item \terme{Outils essentiels} : sélection, recadrage, retouche, peinture, dessin vectoriel, texte, navigation.
    \item \terme{Création, ouverture, enregistrement} : opérations fondamentales.
    \item \terme{Raccourcis clavier} : gain de temps, fluidité, professionnalisme.
    \item \terme{Organisation de l'espace de travail} : espaces prédéfinis, espaces personnalisés, bonnes habitudes.
\end{itemize}

Photoshop est un outil puissant, mais il n'est qu'un outil. La maîtrise de Photoshop ne remplace ni la culture visuelle, ni la créativité, ni la connaissance des principes du design. Elle les prolonge, elle les amplifie. Un designer qui maîtrise Photoshop peut donner vie à ses idées avec précision et efficacité. Un designer qui maîtrise Photoshop sans culture visuelle ne produira que des visuels vides.

Ce chapitre t'a donné les bases nécessaires pour utiliser Photoshop. Mais la maîtrise vient avec la pratique. N'hésite pas à expérimenter, à explorer, à te tromper, à recommencer. Chaque projet est une occasion d'apprendre.

Dans le chapitre suivant, nous aborderons les \terme{documents, images et transformations}. Tu apprendras à créer et gérer des documents, à redimensionner et recadrer des images, à utiliser les repères et les règles, à préparer des documents pour l'écran et pour l'impression. Ces compétences te permettront de travailler efficacement dans Photoshop.

\vspace{2em}

\begin{center}
\textit{Fin du chapitre 6 --- Introduction à Adobe Photoshop.}

\vspace{1em}

\textbf{Prochain chapitre : Chapitre 7 --- Documents, images et transformations.}
\end{center}

